% paper/sections/related_work.tex
% Style: [Author et al., year] showed/proposed... 1-2 sentences each.
%        Gap introduced by "However," or "Yet,".

\section{Related Work}
\label{sec:related}

Revenue maximization in social networks builds on three lines of research:
influence maximization, mechanism design for network markets, and deep RL on graphs.

\paragraph{Influence maximization.}
\citet{kempe2003maximizing} formulated seed-set selection under the Independent
Cascade and Linear Threshold models, proving a $(1 - 1/e)$-approximation via
greedy hill climbing.
Subsequent work scaled this to large graphs through lazy evaluations, sketching,
and diffusion simulations~\citep{goyal2011celf++,tang2014influence,tang2015influence}.
\citet{dai2017learning} introduced S2V-DQN, where a structure2vec GNN produces
node embeddings and a DQN selects seeds; this line was extended to
ToupleGDD~\citep{touplegdd}, which uses graph attention and policy
gradients to achieve state-of-the-art spread on large networks.
Fair variants of influence maximization~\citep{fish2019gaps,tsang2019group}
restrict seed selection to satisfy demographic constraints;
\citet{khajehnejad2022crosswalk} reweight random-walk probabilities to equalize
cross-group diffusion.
However, all of these methods maximize \emph{spread}, not \emph{revenue}, and
assume zero production cost and no pricing structure.

\paragraph{Revenue maximization in networks.}
\citet{babaei2013revenue} showed that a seller can collect strictly more revenue by
offering personalized discounts rather than giving items free, because a buyer's
willingness-to-pay increases as more of her contacts purchase.
Their algorithm selects seeds greedily by degree and assigns prices via a hand-crafted
degree rule (the $\mu$-discount), keeping selection and pricing independent.
Subsequent work by \citet{singer2012win} and \citet{hartline2008optimal} studied
pricing in cascade models under different information assumptions, but focused on
optimal single-item auctions rather than sequential joint policies.
\citet{kempe2014maximizing} extended the revenue framework to general buyer
valuations, but again with decoupled seed selection and pricing.
We are the first to learn a joint policy end-to-end.

\paragraph{Graph neural networks for combinatorial optimization.}
\citet{khalil2017learning} used GNNs with policy gradients to solve graph
combinatorial problems such as Minimum Vertex Cover and Traveling Salesman.
\citet{ding2020erdos} combined GNNs with REINFORCE for the Maximum Cut problem.
Our work applies the same graph RL paradigm to the continuous-discrete joint action
$(v^*, d_t)$ of seed selection and personalized discount, with the additional
complexity of a valuation model that depends on the current seed set.

\paragraph{Sequential pricing and mechanism design.}
The posted-price mechanism literature~\citep{blumrosen2011posted,roughgarden2012supply}
studies how a seller sets prices for buyers arriving in sequence; revenue maximization
in this setting is studied under various buyer-arrival
models~\citep{chawla2010multi}.
Budget-limited influence maximization has been studied in the context of
advertising~\citep{nguyen2013budget,zhang2016maximizing}, but the budget there
limits the \emph{number} of seeds rather than the capital available for production.
Our formulation introduces a \emph{replenishable} budget that grows with accepted
sales, which to our knowledge has not appeared in prior work.

\paragraph{LSTM and Transformer for sequential optimization.}
Sequence models have been applied to combinatorial routing via pointer
networks~\citep{vinyals2015pointer} and to vehicle routing with
transformers~\citep{kool2019attention}.
\citet{chen2019learning} combined GNNs with LSTMs for local search on combinatorial
problems.
We use an LSTM to carry the episode history (past offers, acceptances, running revenue)
across steps, providing the policy with a compressed summary of how the network
responds to pricing, information that a stateless GNN cannot access.
